\section{What the executable has to prove}
\label{l3:contract}
\noindent\textbf{Aim.} Read the maintained implementation as a sequence of
mathematical assertions. A suggested 90-minute session devotes 15 minutes
to the proof contract, 35 to the entry verifier, 20 to the Fourier verifier,
and 20 to reproduction and results. The function descriptions also serve
as a reference for a longer practical session with the source open.

For the circumradius-one regular polygon, let
\[
 -\Delta u=1,\qquad u|_{\partial P_n}=0,\qquad
 J=\frac12\int_{P_n}u,\qquad \F=J/|P_n|^2.
\]
The conclusion sought is that the vertex Hessian of $\F$ has exactly
$2n-4$ strictly negative eigenvalues and the four zero eigenvalues imposed
by two translations, rotation, and dilation. Lecture 1 explains why this
implies strict local maximality modulo similarities. Lecture 2 supplies
the error bounds for the individual Hessian entries.

There are two different finite computations. The first encloses exact
PDE-dependent Hessian entries. The second converts those enclosures into
eigenvalue signs. The following diagram is also the order in which to read
the implementation.
\begin{center}
\begin{tikzpicture}[node distance=6mm,
 box/.style={draw,rounded corners,align=center,text width=11.2cm,
 inner sep=6pt,font=\small},>=Stealth]
\node[box] (a) {FreeFEM + rotations: candidate state, material derivatives,\\
 second lifting, and four scalar flux potentials};
\node[box,below=of a] (b) {\texttt{second\_lifting\_cert}: exact regular fan and directions,\\
 exact element integration, residual and flux error bounds};
\node[box,below=of b] (c) {Arb containment checks: every entry enclosure is inside\\
 the reference ball that will actually be passed to the mode verifier};
\node[box,below=of c] (d) {\texttt{mode\_cert}: all $n$ Fourier blocks, exact symmetry zeros,\\
 $2n-4$ negative eigenvalues and four known zeros};
\draw[->] (a)--(b);\draw[->] (b)--(c);\draw[->] (c)--(d);
\end{tikzpicture}
\end{center}

FreeFEM and NumPy propose finite element functions. Their linear solves,
floating geometry, and FFTs need not be exact: the C verifier checks the
geometry and directions, then bounds the errors of the supplied functions.
The mathematical identities, the implementation of the two verifiers, and
FLINT/Arb's arithmetic remain part of the trusted argument. This is a
computer-assisted proof, rather than a proof object checked by a formal
proof assistant.

The current source root is \path{MaxTorsionValidation/}. The accompanying
\path{source_index.md} gives function locations, and \path{sources.json}
records the hashes of the files used for these notes. Descriptions below
refer to that snapshot, so a reader can distinguish it from sources saved
inside an older computation archive.

\section{The finite element data entering one entry check}
Write $\vartheta=2\pi/n$. Subdivide each coarse fan triangle
\[
 T_j=\operatorname{conv}\{0,(\cos j\vartheta,\sin j\vartheta),
                  (\cos((j+1)\vartheta),\sin((j+1)\vartheta))\}
\]
uniformly with $m$ segments on each side. Then
\[
 \#\text{triangles}=nm^2,\qquad
 \#\text{nodes}=1+\frac{nm(m+1)}2,\qquad
 \#\text{boundary nodes}=nm.
\]
Every fine triangle is contained in one coarse sector. Thus the gradients
of the fan extensions and the pullback coefficients are constant on each
fine triangle. This fitted structure is used both in the regularity-based
rate discussion and in the exact coefficient reconstruction.

An export starts with
\begin{lstlisting}
TORSION_SECOND_V2 n m k qa qp ra rp nv nt
\end{lstlisting}
Here $k$ is the Fourier index; \code{qa}, \code{ra} select radial ($0$) or
tangential ($1$) directions, and \code{qp}, \code{rp} select cosine ($0$)
or sine ($1$). The final two integers are the node and triangle counts.
Next come the node coordinates, the three node indices of each triangle,
and twelve nodal arrays, in this order:
\begin{lstlisting}
qx qy rx ry u uq ur zh psi0 psiq psir psiz
\end{lstlisting}
The first four arrays are components of the fan extensions
$\theta_q,\theta_r$. The next four specify conforming candidates for the
state, two first material derivatives, and the second lifting. The final
four specify continuous piecewise affine scalar potentials whose curls
are used in the equilibrated fluxes. An array named \code{u} is therefore
a computed candidate, not the exact torsion solution.

For a nonzero index other than the even Nyquist index, the normalized
real directions are
\begin{align*}
 r^c_{k,j}&=\sqrt{2/n}\cos(kj\vartheta)
                    (\cos j\vartheta,\sin j\vartheta),
                    &&\text{radial, cosine},\\
 t^c_{k,j}&=\sqrt{2/n}\cos(kj\vartheta)
                    (-\sin j\vartheta,\cos j\vartheta),
                    &&\text{tangential, cosine},\\
 t^s_{k,j}&=\sqrt{2/n}\sin(kj\vartheta)
                    (-\sin j\vartheta,\cos j\vartheta),
                    &&\text{tangential, sine}.
\end{align*}
Using the Fourier convention of Lecture 1, the four requested entries are
the following real Hessian pairings:
\begin{center}
\begin{tabular}{@{}llll@{}}
\toprule
Entry name & $q$ & $r$ & Exact value of $D^2\F[q,r]$\\\midrule
\code{rr} & $r_k^c$ & $r_k^c$ & $\alpha_k$\\
\code{tt} & $t_k^c$ & $t_k^c$ & $\beta_k$\\
\code{rt_re} & $r_k^c$ & $t_k^c$ & $\operatorname{Re}(B_k)_{12}=0$\\
\code{rt_im} & $r_k^c$ & $t_k^s$ & $\operatorname{Im}(B_k)_{12}=\gamma_k$\\
\bottomrule
\end{tabular}
\end{center}
For $n$ even and $k=n/2$, replace
$\sqrt{2/n}$ by $1/\sqrt n$ and omit the zero sine direction. Only the two
diagonal pairings are then needed.

\section{Arb arithmetic: enclosures rather than rounded answers}
An Arb real ball represents a midpoint with a rigorous absolute error
radius. Arithmetic propagates enclosures; increasing the precision
reduces rounding error but does not remove finite element error.
The maintained driver uses 192 bits in both verification stages.
The relevant documented operations are: \code{arb_set_str} for decimal
input, \code{arb_add_error} for enlarging a radius,
\code{arb_contains} for containment, and \code{arb_is_negative} for strict
negativity of the whole ball \cite{FLINTArbDocs}.

In particular, a negative midpoint is insufficient. If an enclosure is
$[-3\cdot10^{-5},10^{-6}]$, its sign is unresolved. Similarly, printed
digits are not the internal ball. The driver may use printed output to
propose a convenient radius, but the entry verifier must subsequently
prove that the exact reference ball contains its internal enclosure.

The code uses \code{arb_sqrtpos} only for expressions mathematically
known to be nonnegative, such as sums of squares. Intersecting their
enclosure with the nonnegative half-line is justified by that identity;
it would not justify taking a square root of an arbitrary uncertain input.
The input readers reject nonfinite numbers before they are used in the
certificate.

\section{The entry verifier, function by function}
\label{l3:entryfunctions}
Open \path{MaxTorsionValidation/flint/second_lifting_cert.c}.
The program does not solve a large interval linear system. It integrates
the supplied piecewise affine fields and evaluates the error inequalities
of Lecture 2. Its functions fall into four short groups, followed by the
main calculation.

\subsection{Reading a finite conforming mesh}
\paragraph{\code{die}.}
Print an explanation and exit with status $2$. A failed input assertion
is an error, not a certificate with a large radius.

\paragraph{\code{cmp_edge} and \code{cmp_tri_key}.}
Compare sorted vertex pairs and triples lexicographically. Sorting makes
edge multiplicities and duplicate triangles detectable independently of
their order in the exported file.

\paragraph{\code{dsu_root} and \code{dsu_join}.}
Maintain connected components using a disjoint-set data structure, with
path compression. They supply the connectivity part of the mesh check.

\paragraph{\code{certify_mesh_topology}.}
Check valid and distinct triangle vertices, reject duplicate triangles,
require all nodes to be used and connected, count edges with one or two
incident triangles, and check the boundary degrees, boundary count, and
Euler relation. In strict mode the node and triangle counts must equal
the regular-fan formulas above. The function returns the boundary mask.
These combinatorial conditions are necessary; the later equality with
the complete canonical triangulation establishes the particular mesh
required by the proof.

\paragraph{\code{read_ball} and \code{read_field}.}
Read a finite decimal as an Arb ball and optionally enlarge it by a
specified radius; repeat for a nodal array. Geometry and direction
candidates are initially given radii $10^{-13}$. PDE and potential
coefficients are read with no additional modelling radius: the decimal
coefficients define the candidate finite element function, with conversion
rounding enclosed by Arb. Truncation and nonfinite data are rejected.

\paragraph{\code{assemble_triangle_geometry}.}
Form the signed determinant of the two triangle edge vectors, prove it
does not contain zero, and compute the area and the gradients of the three
barycentric basis functions. The signed determinant is used for gradients;
its absolute value is used for area. Geometry is assembled again after
the exact regular coordinates have replaced the candidate coordinates.

\subsection{Proving the exact regular fan and the directions}
\paragraph{\code{canonical_gid}.}
Give each lattice node of the regular fan a single integer index. The
centre has one index; nodes on shared rays have the same index when
described from either neighboring sector. This is how the program avoids
counting a shared ray twice.

\paragraph{\code{canonical_tri_key}.}
Sort the three canonical node indices of a triangle. This removes an
irrelevant orientation or input-order choice when comparing triangulations.

\paragraph{\code{validate_regular_inputs}.}
Reconstruct the exact node set
\[
 x=\frac{s}{m}(\cos j\vartheta,\sin j\vartheta)
   +\frac{t}{m}(\cos((j+1)\vartheta),\sin((j+1)\vartheta)),
 \qquad s,t\geq0,\quad s+t\leq m.
\]
The integer pair $(s,t)$ is a lattice address, not a new angle notation.
Floating calculations nominate an address; Arb then proves that the
candidate coordinate balls contain the trigonometric coordinates of that
address. The assignment must be a bijection. The boundary mask must agree
with the canonical boundary, and the sorted list of triangles must equal
the full canonical list, not merely have the correct cardinality.

The function also reconstructs the normalized Fourier directions at each
node and proves their containment in the exported direction balls.
It finally replaces geometry and direction arrays by these analytic
enclosures. Each triangle has exact area $\sin\vartheta/(2m^2)$ and the
total area is $n\sin\vartheta/2$. Consequently the $10^{-13}$ input radii
serve to identify the intended exact data; they do not certify an arbitrary
perturbation of the polygon or enter as a shape uncertainty.

\subsection{Exact element integration and coefficient bounds}
\paragraph{\code{gradient}.}
For a nodal scalar field $v$, compute
$\nabla v|_T=\sum_{i=0}^2v_i\nabla\lambda_i$. It is constant on a
triangle, so no numerical quadrature is required.

\paragraph{\code{affine_integral}.}
Use $\int_Tv=|T|(v_0+v_1+v_2)/3$.

\paragraph{\code{add_affine_square}.}
Add the exact integral
\begin{equation}
 \int_Tv^2=\frac{|T|}{6}
 (v_0^2+v_1^2+v_2^2+v_0v_1+v_0v_2+v_1v_2).
 \label{l3:mass}
\end{equation}
This is the usual consistent $P_1$ mass matrix. As a quick check,
$v_0=v_1=v_2=1$ gives $|T|$. Applied separately to the two components of
an affine flux mismatch, it gives its squared $L^2$ norm exactly, with
only interval arithmetic rounding to enclose.

\paragraph{\code{symmetric_norm}.}
For a real symmetric matrix with entries $x,y,z$, compute its operator norm
from
\[
 \left\|\begin{pmatrix}x&y\\y&z\end{pmatrix}\right\|_2
   =\frac{|x+z|+\sqrt{(x-z)^2+4y^2}}2.
\]
Taking a maximum over triangles supplies the constants
$C_q=\|\mathbb A_q\|_{L^\infty}$, $C_r$, and
$C_{qr}=\|\mathbb A_{qr}\|_{L^\infty}$ used in Lecture 2.

\paragraph{\code{coefficients}.}
Differentiate the planar pullback coefficient
\[
 \mathbb A(s,t)=\det(\Id+sD\theta_q+tD\theta_r)
 (\Id+sD\theta_q+tD\theta_r)^{-1}
 (\Id+sD\theta_q+tD\theta_r)^{-T}.
\]
The function returns $\mathbb A_q,\mathbb A_r,\mathbb A_{qr}$ and
the mixed determinant derivative
\[
 (\operatorname{div}\theta_q)(\operatorname{div}\theta_r)
                          -\tr(D\theta_qD\theta_r).
\]
The first matrix derivative is
$\mathbb A_q=(\operatorname{div}\theta_q)\Id-D\theta_q-D\theta_q^T$.
To read the second derivative without an extra chain of aliases, expand it
as
\begin{align*}
 \mathbb A_{qr}={}&
 \bigl((\operatorname{div}\theta_q)(\operatorname{div}\theta_r)
                       -\tr(D\theta_qD\theta_r)\bigr)\Id\\
 &-(\operatorname{div}\theta_q)(D\theta_r+D\theta_r^T)
  -(\operatorname{div}\theta_r)(D\theta_q+D\theta_q^T)\\
 &+D\theta_qD\theta_r+D\theta_rD\theta_q
  +D\theta_q^TD\theta_r^T+D\theta_r^TD\theta_q^T\\
 &+D\theta_qD\theta_r^T+D\theta_rD\theta_q^T.
\end{align*}
The source stores each symmetric matrix in three entries. Its temporary
four-entry arrays named \code{Q} and \code{R} are simply the row-major
Jacobians $D\theta_q,D\theta_r$; no additional mathematical quantities
are needed to understand them.

\paragraph{\code{state_error}.}
Compute
\[
 \Phi_0=\|-x/2+\curl\widetilde\psi_0-\nabla\widetilde u\|_{L^2(P_n)},
 \qquad \curl\psi=(\partial_y\psi,-\partial_x\psi).
\]
The flux has divergence $-1$. Continuity of the scalar potential ensures
the required normal continuity of its curl. Thus $\Phi_0$ bounds the
state energy error even if the auxiliary fit for $\widetilde\psi_0$ is
inexact.

\subsection{Residuals, reference values, and the main calculation}
\paragraph{\code{interior_residual_norm}.}
Take the Euclidean norm of a residual array on the interior nodes only.
The basis of $V_h\subset H^1_0(P_n)$ has no boundary degrees of freedom.
Including the boundary rows of the unconstrained stiffness matrix would
therefore compute the wrong algebraic residual for this problem.

\paragraph{\code{read_mode_reference}.}
Read the requested component of a \code{MODE_DATA} record: \code{rr},
\code{tt}, \code{rt_re}, or \code{rt_im}. These correspond respectively
to the two diagonal entries and the real and imaginary parts of the
upper-right entry. The entry pass checks the requested index and its
conjugate; the final mode pass checks uniqueness and completeness of all
Fourier records.

\paragraph{\code{main}: assertions before the triangle loop.}
Require exactly one of \code{--check-regular-inputs} and
\code{--diagnostic-nonstrict}. Reference certification requires the strict
choice, the reference file, the entry name, and the proposed maximum
radius. Validate dimensions, metadata, and radii; read the export; reject
trailing data. Set the four Dirichlet candidate fields to exact zero on
the certified boundary. The scalar flux potentials have no Dirichlet
condition. Then validate and reconstruct the exact regular inputs.

\paragraph{\code{main}: the triangle loop.}
Accumulate the area, $\int\widetilde u$, the second area variation, the
raw Hessian centre, the three coefficient norm bounds, the four residual
vectors, the node patch areas, and the three remaining squared flux
mismatches. The state mismatch was already evaluated by
\code{state_error}. Each calculation uses constant gradients and
\eqref{l3:mass}; there is no quadrature remainder to estimate.

The mathematical residual convention is
$\rho=\text{right side}-\text{matrix}\times\text{candidate}$.
The C arrays \code{res0}, \code{resq}, \code{resr}, \code{resz}
store its \emph{negative}. Norms do not change, but the sign matters when
correcting the centre.

\paragraph{\code{main}: algebraic errors.}
The polygon lies in $(-1,1)^2$, so domain monotonicity gives
$\lambda_1(P_n)\geq\pi^2/2$. The local mass matrix gives, for the
interior-node stiffness matrix $K_h$,
\[
 K_h\succeq\alpha\Id,\qquad
 \alpha=\frac{\pi^2}{24}
       \min_{i\ \mathrm{interior}}\sum_{T\ni i}|T|>0.
\]
Thus the dual norm of a residual restricted to $V_h$ is bounded by its
Euclidean coefficient norm divided by $\sqrt\alpha$. The first-derivative
algebraic bounds also contain the state error multiplied by $C_q$ or
$C_r$, because the right sides depend on the state candidate.

The corrected centres are those derived in Lecture 2:
\[
 \widetilde C_{\rm corr}
  =\widetilde C+\rho_r(\widetilde u_q)+\rho_0(\widetilde z),
 \qquad
 \widetilde J_{\rm corr}
  =\int_P\widetilde u-\tfrac12\int_P|\nabla\widetilde u|^2.
\]
In the C loop, adding $\rho_r(\widetilde u_q)$ means subtracting
\code{resr} dotted with \code{uq}. The remaining algebraic error in the
Hessian centre is quadratic in the relevant energy errors.

\paragraph{\code{main}: the final entry enclosure.}
Let $\widehat\delta_0,\widehat\delta_q,\widehat\delta_r$ and
$\widehat\eta_{qr}$ denote the computed upper bounds for the state,
first-derivative, and second-lifting Galerkin errors, and let
$E_{\rm alg}$ bound the corrected discrete-centre error. These are the
quantities of Lecture 2, not new error definitions. For the nonzero
Fourier pairings the first area variations vanish. The program evaluates
\begin{align*}
 \widetilde F_{qr}
  &=\frac{\widetilde C_{\rm corr}}{A^2}
        -\frac{2\widetilde J_{\rm corr}A_{qr}}{A^3},\\
 R_F&=\frac{\widehat\delta_q\widehat\delta_r
       +\tfrac12 C_{qr}\widehat\delta_0^2
       +\widehat\delta_0\widehat\eta_{qr}+E_{\rm alg}}{A^2}
       +\frac{|A_{qr}|}{A^3}
                    (\widehat\delta_0^2+\epsilon_0^2).
\end{align*}
Here $A=|P_n|$, $A_{qr}=D^2A[q,r]$, and $\epsilon_0$ bounds the state
algebraic energy error. It enlarges the ball for $\widetilde F_{qr}$ by
$R_F$. Its output separates \code{F_hessian_continuous_error<=} from
\code{F_hessian_algebraic_error<=}, making the source of a large radius
visible.

\paragraph{\code{main}: the handoff assertion.}
If a reference is supplied, reconstruct its ball using the proposed common
mode radius and prove that it contains the complete entry enclosure.
Repeat at the conjugate mode $n-k$; negate the imaginary entry there.
Only successful containment prints
\begin{lstlisting}
ENTRY_CERTIFIED k=... conjugate_k=... entry=... max_radius=...
\end{lstlisting}
The driver checks that all four metadata fields match its request. This
second pass closes the gap between a rigorously computed internal ball
and the finite decimal data actually sent to the eigenvalue verifier.

\section{The Fourier verifier, function by function}
Open \path{MaxTorsionValidation/flint/mode_cert.c}. Its input consists of
\begin{lstlisting}
MODE_DATA k a d Re(b) Im(b)
\end{lstlisting}
representing $\begin{pmatrix}a&b\\\overline b&d\end{pmatrix}$, and a
radius file with one \code{MODE_RADIUS k radius} line for every index.
The letters in this file format mean the two diagonal entries and the
complex upper-right entry; in Lecture 1 these are $\alpha_k,\beta_k$,
and $i\gamma_k$.

\paragraph{\code{usage}.}
Describe the accepted options and exit with status $2$. A global radius
and a per-mode radius file cannot be used together.

\paragraph{\code{print_ball}.}
Print an enclosure for inspection. The displayed precision does not
replace the precision used in the sign predicates.

\paragraph{\code{main}: input validation.}
Read finite real components and nonnegative radii. With the expected
number of symbols specified, require exactly one mode and one radius
for every $k=0,\ldots,n-1$. Reject missing, duplicate, or out-of-range
indices. These checks prevent a repeated negative mode from standing in
for an untested one.

\paragraph{\code{main}: analytic symmetry information.}
With \code{--exact-regular-similarities}, replace $B_0$ by the exact zero
matrix. At $k=1,n-1$, use the exact translation eigenvalue zero and the
trace for the other eigenvalue. At the even Nyquist index, set the
off-diagonal entries to zero. These are analytic identities from Lecture 1;
they are applicable because the first stage proved that the data correspond
to the regular polygon and its normalized Fourier directions.

\paragraph{\code{main}: eigenvalues and sign counts.}
For the remaining blocks evaluate
\[
 \lambda_\pm=
 \frac{a+d\pm\sqrt{(a-d)^2+4(\operatorname{Re}b)^2
                                  +4(\operatorname{Im}b)^2}}2.
\]
Count an eigenvalue as negative only when its entire Arb ball is strictly
negative. The internal word \code{unresolved} includes the four exact
zeros. In a successful regular-polygon run these are its only members.
An interval proved positive is not counted as unresolved; it also cannot
satisfy the required negative count.

The full driver requires all three counts
\[
 \text{symbols}=n,\qquad \text{negative}=2n-4,
                  \qquad \text{unresolved}=4.
\]
The program exits $0$ when the requested counts hold, $3$ when they fail,
and $2$ on an input error. Its printed \code{MODE_CERTIFIED} marker means
the \emph{requested} counts were satisfied. A standalone invocation that
requests weaker counts does not prove local maximality. Nor does this
program prove the PDE error radii: those must come from the entry checks.

\section{How the rotation-based driver obtains the candidates}
\label{l3:driver}
The optimization concerns the first derivatives. At vertex $0$, the local
radial and tangential directions are $(1,0)$ and $(0,1)$. Solve the torsion
problem and these two material problems. Rotational symmetry gives the
other local vertex derivatives by composing these scalar functions with
the inverse rotation. Weighted sums give the Fourier derivatives.
Thus only three primal state/first-derivative solves are needed per polygon.
The present error estimator still solves a second lifting for each real
Hessian pairing, and it computes auxiliary scalar flux fits. Counting
those as well would give more than three linear solves.

\subsection{The FreeFEM problems}
In \path{freefem/second_variation_lifting_rotations.edp}, the named problem
\code{state} is the torsion Galerkin equation, \code{firstQ} and
\code{firstR} are the two differentiated Galerkin equations, and
\code{lifting} is the discrete second lifting of Lecture 2.
The \code{-make-basis 1} branch solves the first three once at vertex $0$.
The \code{-cached-fields} branch reads the state and the requested first
derivatives from the cache and solves only \code{lifting}.

The named problem \code{fitCurl} minimizes a regularized flux mismatch over
continuous piecewise affine scalar potentials. The added term is
$10^{-12}\int_P\psi_h^2/2$ in the minimized functional: it makes the fit
coercive and fixes the additive constant, while also slightly changing the
nonconstant part of the minimizer. Three such fits accompany the initial basis;
one more accompanies each second lifting. The argument for equilibrium
uses the form of the flux, not successful convergence of this optimization.
The solver tolerance affects sharpness and speed, while the C verifier
evaluates the actual supplied potential.

The separate \path{freefem/torsion_hessian_rotations.edp} assembles a
floating Hessian from the same three-solve principle. It is useful for
inspection and comparisons. Its named problems \code{torsion},
\code{materialX}, and \code{materialY} solve the state and the two
coordinate material equations at vertex zero. The driver option \code{--entry-centers}
avoids needing this additional floating Hessian by taking reference
centres directly from the entry estimates.

\subsection{\code{rotation_cache.py}}
\paragraph{\code{prepare}.}
Read the basis export, check dimensions and finiteness, identify lattice
addresses, construct the node permutation for rotation, and check its
cycles and coordinate consistency. Shared-ray addresses have one canonical
key. Apply the permutation to the two vertex-zero scalar solutions and
their flux potentials, then form Fourier combinations with NumPy's FFT.
For NumPy's negative-exponent convention, cosine weights come from the
real part and positive sine weights from the negative imaginary part.
The normalization is $\sqrt{n/2}$, or $\sqrt n$ for the even Nyquist mode,
when dividing an unnormalized Fourier sum.

The nested function \code{direction} selects the radial/tangential and
cosine/sine combination for one requested field. The cache is only a
candidate-generation device: exact geometric and directional checks,
and residual bounds for the resulting fields, are repeated in C.

\subsection{\code{certificate_build.py}}
\paragraph{\code{build_verifiers}.}
Compile \code{second_lifting_cert} and \code{mode_cert} from the source
copies saved inside the new archive, using a forced build. Return their
private executable paths to the driver and record their hashes. A small
diagnostic records the compiler and FLINT header/runtime versions. Failure
of that diagnostic is recorded as unavailable metadata; the two verifier
builds themselves must succeed. The run therefore does not silently use
a pre-existing executable from a different source version.

\subsection{\code{certify_with_rotations.py}}
\paragraph{\code{run}.}
Execute a subprocess, save its output in a designated log, measure elapsed
time, and raise an error for a nonzero exit status. The final mode step is
handled separately so that an inconclusive certificate can be recorded
normally rather than mistaken for a crashed program.

\paragraph{\code{ball}.}
Read the printed midpoint and radius for a named output using decimal
arithmetic. Its result is used to propose a radius. It has no independent
authority to prove containment.

\paragraph{\code{main}: setup and candidate phase.}
Parse the range of $n$, $m$, worker count, CPU affinity limit, flux-fit
tolerance, and a new archive path. Limit numerical-library threads to one
per process, set the chosen affinity, and run with reduced scheduling
priority. Save a source snapshot and manifest, build the private verifiers,
and create the three-solve rotation cache for each polygon.

The nested function \code{evaluate} generates one second-lifting export,
checks its header against the requested pairing, and runs the entry
verifier in strict mode. From the displayed centre and error ball it
proposes a slightly enlarged decimal radius covering both $k$ and $n-k$.
Four pairings are used for each representative nonzero mode, except the
even Nyquist mode, which needs two.

\paragraph{\code{main}: reference and verification phase.}
Once all centres exist, write \code{modes.txt} once and retain it unchanged
during verification. Take the maximum proposed entry radius for each mode
and its conjugate. The nested function \code{verify} reruns the C entry
verifier with this exact reference file and radius, checks the complete
\code{ENTRY_CERTIFIED} marker, and compresses the checked export.
After all entries pass, run \code{mode_cert} at 192 bits with the required
$n,2n-4,4$ counts. Save \code{RESULT.json}, append \code{results.jsonl},
and stop on an inconclusive polygon if \code{--stop-on-failure} is set.
Finally save archive hashes and a \code{FINISHED} marker.

\paragraph{Reading a run's status.}
The Python driver can exit normally after recording
\code{NOT_CERTIFIED}; this is an expected numerical outcome.
\code{FINISHED} means that the requested sweep or its stopping rule has
finished. To claim success for a particular $n$, inspect its
\code{RESULT.json}: it must say \code{status: CERTIFIED} and
\code{mode_exit_code: 0}, with the corresponding entry logs and final
counts present. A shell exit code alone is insufficient.

\section{Other validation utilities and their precise roles}
The two C programs above are the maintained complete certification path.
The repository also contains smaller verifiers and diagnostic utilities.
Their output should be read according to the narrower assertion each
one actually establishes.

\subsection{\code{residual_cert.c}}
\code{read_arb} parses a scalar. Its \code{main} reads a dense matrix,
right side, candidate vector, and a supplied positive lower bound
$\alpha$; it evaluates the residual and the resulting algebraic error
bound. The assumption that the matrix is symmetric positive definite and
that $\alpha\leq\lambda_{\min}$ is supplied by the caller. Positivity of
the number $\alpha$ alone does not prove that assumption. This utility
does not perform the regular-fan and boundary checks of the main verifier.

\subsection{\code{second_variation_cert.c}}
\code{die} reports invalid input; \code{set_nonnegative} parses and checks
nonnegative bounds. The \code{main} function reads complete records of
two direction indices and five bounds, and evaluates
\[
 \delta_q\delta_r+\tfrac12 C_{qr}\delta_0^2+\delta_0\eta_{qr}.
\]
It checks the arithmetic combination, assuming that the five supplied
bounds have already been proved. It does not establish their PDE meaning.

\subsection{\code{majorant_cert.c}}
\begin{longtable}{@{}p{.34\textwidth}p{.62\textwidth}@{}}
\toprule Function & Role\\\midrule\endhead
\code{die} & Reject an invalid export.\\
\code{read_arb}, \code{read_field} & Read scalar balls and arrays.\\
\code{grad_field} & Compute a $P_1$ gradient on a triangle.\\
\code{add_affine_square} & Integrate a squared affine mismatch.\\
\code{energy_norm} & Integrate a field's squared gradient and take its norm.\\
\code{state_majorant} & Bound the state error by its equilibrated flux mismatch.\\
\code{material_majorant} & Bound a first material error, including propagation of the state error.\\
\code{print_ball} & Display an enclosure.\\
\code{main} & Read candidate data and combine direct state/material majorants into diagnostic bounds.\\
\bottomrule
\end{longtable}
This program is useful for understanding individual majorants and direct
Gram estimates. It does not establish the complete canonical-fan input
contract used by \code{second_lifting_cert}; its output alone is not the
maintained local-maximality certificate.

\subsection{Drivers, summaries, and checks}
\path{freefem/test_extended_certificates.py} follows the same two-stage
certificate structure with per-entry candidate solves instead of a shared
rotation cache. Its \code{run}, \code{ball}, nested \code{evaluate} and
\code{verify}, and \code{main} have the corresponding roles described
above. The original candidate-generating FreeFEM files remain available
for comparisons.

In \path{freefem/summarize_extended_certificates.py}, \code{audit} reads an
archive's metadata, source hashes, entry markers, radii, and mode output;
\code{main} assembles the requested result tables. Table upper endpoints
are rounded outward using decimal arithmetic. This is an audit of records,
not a fresh replay of every finite element export.

The input-rejection tests in \path{flint/test_certificate_inputs.py} and
the rotation comparisons in
\path{freefem/test_torsion_hessian_rotations.py} exercise implementation
properties. Neither a successful comparison with floating output nor a
successful test suite substitutes for the entry enclosures and sign counts
for a particular polygon.

\section{Running a certificate yourself}
The commands below are run from the repository root. They require
FreeFEM, a C compiler with FLINT/Arb development files, Python, and NumPy;
the maintained installation instructions are in
\path{MaxTorsionValidation/README.md}. The archive directory must not
already exist, so each run preserves its own evidence.

For a small first run:
\begin{lstlisting}[language=bash]
python3 MaxTorsionValidation/freefem/certify_with_rotations.py \
  5 5 --m 32 --entry-centers --jobs 6 --cpu-cores 6 \
  --stop-on-failure --archive /tmp/torsion_n5_m32_course
\end{lstlisting}
For the largest successfully archived case:
\begin{lstlisting}[language=bash]
python3 MaxTorsionValidation/freefem/certify_with_rotations.py \
  25 25 --m 128 --entry-centers --jobs 6 --cpu-cores 6 \
  --flux-eps 1e-13 --stop-on-failure \
  --archive /tmp/torsion_n25_m128_course
\end{lstlisting}
To test a range, replace the two positional integers by its first and
last $n$. The stopping option halts after the first failure to certify.
For this project's six-physical-core machine the shown settings use six
workers. More generally, \code{--cpu-cores} selects from OS CPU affinity
indices; it does not discover physical-core topology on another machine.

Inspect a completed run with, for example,
\begin{lstlisting}[language=bash]
cat /tmp/torsion_n25_m128_course/n25/RESULT.json
tail -n 2 /tmp/torsion_n25_m128_course/n25/mode_cert.log
\end{lstlisting}
The expected final counts for $n=25$ are $25$ symbols, $46$ negative
eigenvalues, and four exact similarity zeros. The archive also contains
the reference centres, per-mode radii, compressed entry exports, entry
logs, source snapshot, build metadata, and hashes.

An independent replay of a single entry uses the archived verifier and
the decompressed export, with the same \code{--prec 192}, geometry and
direction radii, strict input option, reference file, entry name, and
\code{--max-radius} as its original command. Running only the final
\code{mode_cert} is a useful replay of the small spectral calculation,
but it does not replay the entry PDE bounds.

\section{What the archived computations establish}
\label{l3:results}
The following tables transcribe the manuscript and the archived table
data. No new certification sweep was run to write these notes. All
eigenvalue values refer to $D^2(J/A^2)$ in the normalized vertex coordinates
of Lecture 1. The listed upper bound is the largest certified upper
endpoint among eigenvalues other than the four similarity zeros. Negative
upper bounds are rounded toward $+\infty$, so the displayed values remain
safe upper bounds.

\begin{center}
\begin{tabular}{@{}rrrr@{}}
\toprule
$n$ & $m$ & Negative / required & Upper bound\\\midrule
3&32&2/2&$-1.00\,10^{-2}$\\
4&32&4/4&$-6.20\,10^{-3}$\\
5&32&6/6&$-3.10\,10^{-3}$\\
6&32&8/8&$-1.70\,10^{-3}$\\
7&32&10/10&$-9.80\,10^{-4}$\\
8&32&12/12&$-5.90\,10^{-4}$\\
9&32&14/14&$-3.60\,10^{-4}$\\
10&32&16/16&$-2.20\,10^{-4}$\\
\bottomrule
\end{tabular}
\end{center}
The $n=3,4$ cases are also known analytically. The computations extend the
local-maximality result through $n=25$. Four similarity zeros accompany
every successful row; they are known exactly from invariance, not inferred
from small decimal numbers.

% Generated from archived table JSON; no new computations.

\subsection*{The $m=64$ computations}
\begin{center}\small
\begin{tabular}{@{}rrrrl@{}}
\toprule
$n$ & Triangles & Negative / required & Upper bound & Outcome\\\midrule
11 & 45056 & 18 / 18 & $-2.11\,10^{-4}$ & certified\\
12 & 49152 & 20 / 20 & $-1.50\,10^{-4}$ & certified\\
13 & 53248 & 22 / 22 & $-1.08\,10^{-4}$ & certified\\
14 & 57344 & 24 / 24 & $-7.99\,10^{-5}$ & certified\\
15 & 61440 & 26 / 26 & $-5.97\,10^{-5}$ & certified\\
16 & 65536 & 28 / 28 & $-4.51\,10^{-5}$ & certified\\
17 & 69632 & 30 / 30 & $-3.44\,10^{-5}$ & certified\\
18 & 73728 & 32 / 32 & $-2.64\,10^{-5}$ & certified\\
19 & 77824 & 30 / 34 & $1.35\,10^{-5}$ & inconclusive\\
20 & 81920 & 28 / 36 & $1.69\,10^{-4}$ & inconclusive\\
\bottomrule
\end{tabular}\end{center}

\subsection*{The $m=128$ computations}
\begin{center}\small
\begin{tabular}{@{}rrrrrl@{}}
\toprule
$n$ & Triangles & Negative / required & Upper bound & Minutes & Outcome\\\midrule
19 & 311296 & 34 / 34 & $-2.73\,10^{-5}$ & 15.7 & certified\\
20 & 327680 & 36 / 36 & $-2.22\,10^{-5}$ & 18.0 & certified\\
21 & 344064 & 38 / 38 & $-1.81\,10^{-5}$ & 20.8 & certified\\
22 & 360448 & 40 / 40 & $-1.49\,10^{-5}$ & 13.5 & certified\\
23 & 376832 & 42 / 42 & $-1.24\,10^{-5}$ & 15.8 & certified\\
24 & 393216 & 44 / 44 & $-1.03\,10^{-5}$ & 19.4 & certified\\
25 & 409600 & 46 / 46 & $-8.71\,10^{-6}$ & 17.7 & certified\\
\bottomrule
\end{tabular}\end{center}


At $m=64$, the archived checks succeed through $n=18$ and are inconclusive
for $n=19,20$. A positive upper endpoint in those rows does not prove a
positive eigenvalue. It means the enclosure is too wide to establish all
the desired negative signs. Refinement to $m=128$ succeeds for every tested
$n=19,\ldots,25$. These computations neither establish a maximal possible
$n$ at either mesh size nor give a conclusion beyond $n=25$.

\subsection{Time and the meaning of the three-solve optimization}
The recorded $m=128$ run for $n=25$ took $1059.4$ seconds, about
$17.7$ minutes, on the six-physical-core Intel Core i7-9750H machine
(2.60 GHz). It has $409600$ triangles, $206401$ nodes, and $48$ real
entry checks. The full optimized $n=22,\ldots,25$ records total about
$66.4$ minutes; allowing roughly $70$ minutes is a useful estimate on that
machine. The $n=24$ record includes a handoff interval, and the earlier
$n=19,20,21$ runs used a different allocation/workflow, so the timings do
not form a controlled scaling benchmark.

The three initial primal solves do not make the whole proof constant
cost in $n$: there are more modes, more mesh triangles, one second lifting
and one auxiliary flux fit per real pairing, and two Arb entry passes.
This explains why the optimization is substantial while certification
still takes minutes.

\subsection{Where the evidence is stored}
The main archive groups are under \path{MaxTorsionValidation/results/}:
\begin{itemize}
\item \path{review_2026_09_08_certificate_archive}: the $m=32$ chain.
\item \path{extended_m128_fast_n19_onward}: the $m=128$, $n=19,20,21$ records.
\item \path{extended_m128_rotations_n22_n30}: the completed $n=22,23,24$
records; the folder name records an initial range, not results through $30$.
\item \path{extended_m128_rotations_n25}: the separate $n=25$ certificate.
\item \path{extended_m64_table_data.json} and
\path{extended_m128_table_data.json}: the audited table rows and their
per-row log references.
\end{itemize}
The packaging rename from \path{code/} to \path{MaxTorsionValidation/}
preserved historical archives byte for byte. Consequently an old JSON
record can still contain a \path{code/results/...} path; resolve that
prefix to the current source root when locating the file. The packaging
provenance explains this explicitly.

The recorded validity review checked 672 entry-log records, replayed 25
final mode checks, and performed four complete entry replays. Its archive
hash audit covered 2701 files and recorded one historical checksum mismatch
for the compressed $n=21$, $k=6$, radial--radial export. That entry was
freshly recertified at the archived radius; the original checksum record
was retained and the replay documented in
\path{validity_review_20260909/n21-replay-provenance.json}.
This is the precise scope of that review: it was not a fresh recomputation
of all 672 entry exports. Hashes establish correspondence of bytes;
the mathematical conclusion comes from the verification chain applied
to the identified data.

\subsection*{Exercises and audit questions}
\begin{exercise}
An FFT bug changes a cached first derivative. Why can the full chain still
be valid if it reports success?
\end{exercise}
\noindent\emph{Answer.} The cached field is merely a conforming candidate.
The C verifier computes its residual and flux mismatch for the prescribed
exact direction. A poor candidate normally widens the enclosure or causes
failure; it is not assumed to solve the intended PDE.

\begin{exercise}
The entry pass gives an enclosure centred at $c$ with radius $r$. The mode
file stores a decimal $c_0$. Is using radius $r$ around $c_0$ sufficient?
\end{exercise}
\noindent\emph{Answer.} In general no: the difference between $c_0$ and the
entry centre must be covered as well. The second C pass proves containment
of the whole internal enclosure in the ball built from the actual mode
file, and also checks the conjugate entry.

\begin{exercise}
Why are four exact zeros compatible with strict local maximality?
\end{exercise}
\noindent\emph{Answer.} They are precisely the tangent directions to
similarities, along which $J/A^2$ is constant. On a local slice fixing
translation, rotation, and scale, the Hessian is strictly negative definite.

\begin{exercise}
What follows if \code{RESULT.json} says \code{NOT_CERTIFIED} at $m=64$?
\end{exercise}
\noindent\emph{Answer.} This chain at this mesh and these candidates has
not proved the required signs. It supplies no counterexample to local
maximality. The successful $m=128$ checks for $n=19,20$ illustrate the
difference between a failed enclosure test and a false mathematical claim.
